%% kapa: exam=1 round=23 year=2012 period=1 form=A subject=economics subject_ko=경제원론 seq=30 number=70
%% source: Q-Net 감정평가사 기출문제 (cst00309) artlSeq=5052578 file=raw/1008/5052578/2012년 제23회 감정평가사 제1차시험 기출문제(1교시 A형).pdf
%% conversion: kordoc 4.24.51 -> tools/gen1.py (auto); review=proofread-sonnet; flags=none
\begin{question}{70}
실물경기변동이론(real business cycle theory)에 관한 설명으로 옳지 않은 것은?
\begin{choices}
  \item 경기변동의 요인으로 기술 충격의 중요성을 강조한다.
  \item 노동시장은 항상 균형을 이룬다.
  \item 경기변동은 시간에 따른 균형의 변화로 나타난다.
  \item 불경기에도 생산의 효율성은 달성된다.
  \item 생산성은 경기역행적(counter-cyclical)이다.
\end{choices}
\end{question}
\answer{5}
